% !PRINTED 632
\weberpage{632}
%<<CONT
\gesperrt{quantity of the field $\Omega(\varepsilon)$ can be brought in only one way into the normal form
\[
\Phi(\varepsilon) = c_0 + c_1\varepsilon + c_2\varepsilon^2 + \cdots + c_{n-2}\varepsilon^{n-2}
\]
in which the coefficients $c$ are contained in $\Omega$.}

In this theorem we can take for $\Omega$ the field that arises from $R$ by the adjunction of $n$ independent variables $x_0, x_1, x_2 \ldots x_{n-1}$; for if $X$ were reducible in this field, so that
\[
X = X_1 X_2,
\]
where $X_1, X_2$ are integral functions of $\varepsilon$ and rational functions of the $x$, then one could put for the variables $x_0, x_1 \ldots x_{n-1}$ such rational numbers that $X_1, X_2$, without changing their degree with respect to $\varepsilon$, pass into functions in $R$ (\S.\,143), and this contradicts the irreducibility of $X$ in $R$.

This field $\Omega$ has various divisors in which $X$ is likewise irreducible, and which can therefore all be taken for $\Omega$ in Theorem 1. Such divisors are obtained if one fixes any \gesperrt{permutation group} $P$ of the $n$ figures $0, 1, 2 \ldots n-1$ and considers the totality of all rational functions of $x_0, x_1 \ldots x_{n-1}$ in $R$ that \gesperrt{admit} the permutations of this group.

The totality of all these functions evidently forms a field, and thus to every permutation group of the $n$ figures one obtains a definite associated field $\Omega$. The quantities of such a field can be represented rationally by the \gesperrt{symmetric functions of the $x$ and by a function belonging to the group $P$}.

We call the field $\Omega$ so determined \gesperrt{the field belonging to the group $P$}. Every such field can take the place of $\Omega$ in 1. (Cf. \S.\,152.)

We can therefore state the second theorem:

\gesperrt{2. If a rational function $\Phi(x_0, x_1 \ldots x_{n-1}, \varepsilon)$ of the field $\Omega(\varepsilon)$ remains unchanged when the permutations of a group $P$ are applied to the indices of the $x$, then the coefficients $c_0, c_1, c_2 \ldots c_{n-2}$ of the normal form of $\Phi$ also admit the permutations of $P$.}

If $\Phi(\varepsilon)$ is any element of the field $\Omega(\varepsilon)$, one obtains the conjugate quantities by putting for $\varepsilon$ any other
%CONT>>
